\section{Conclusion}
This paper presents an in-depth empirical study of the correctness of plausible generated patches on SWE-bench.
We first find that the validation process of SWE-bench overlooks non-modified test files, which leads to significant performance overestimation. 
Then we extensively investigate the prevalence of plausible patches that exhibit behavioral discrepancies from ground truth patches, the patch difference patterns leading to behavioral discrepancies, and the correctness of plausible patches that exhibit behavioral discrepancies.
The core of our methodology is the novel \appname technique for differential patch testing, which automatically exposes behavioral discrepancies between two patches.
The results call for carefully selecting developer tests for patch validation, checking and filtering out plausible but incorrect patches for more accurate evaluation, and paying more attention to supplementary semantic changes in plausible patches.
Our work will contribute toward better issue-solving tools and benchmarks that address the problem of under-specified specifications.
\appname can provide concrete evidence for invalid functionality in patches through generated tests, enabling more focused and objective patch assessment.
The test generated by \appname can be continuously used to strengthen issue-solving benchmarks with the help of benchmark users.
We envision \appname to be useful for building a sustainable and robust patch validation ecosystem for assessing issue-solving tools.



